\documentclass[11pt,a4paper]{article}
\usepackage[T1]{fontenc}
\usepackage[utf8]{inputenc}
\usepackage{lmodern}
\usepackage[margin=2.45cm,headheight=14pt]{geometry}
\usepackage{amsmath,amssymb,amsthm,mathtools}
\usepackage{microtype,booktabs,longtable,array}
\usepackage{enumitem}
\usepackage{tikz}
\usetikzlibrary{arrows.meta}
\usepackage{xcolor}
\usepackage{fancyhdr,placeins}
\usepackage[colorlinks=true,linkcolor=blue!45!black,citecolor=blue!45!black,urlcolor=blue!45!black]{hyperref}
\hypersetup{pdftitle={Nonintegral Tropical Degree from Two Horizontal Segments},pdfauthor={AI-assisted research draft},pdfsubject={An exact tropical Hilbert function and proposed counterexamples to two questions of Grigoriev}}
\pagestyle{fancy}
\fancyhf{}
\fancyhead[L]{\small Nonintegral tropical degree}
\fancyhead[R]{\small Research draft, 20 September 2026}
\fancyfoot[C]{\thepage}
\renewcommand{\headrulewidth}{0.3pt}
\setlength{\parskip}{0.28em}
\setlength{\parindent}{1.25em}
\setlength{\emergencystretch}{2em}
\setlist{itemsep=0.25em,topsep=0.35em}
\numberwithin{equation}{section}
\newtheorem{theorem}{Theorem}[section]
\newtheorem{proposition}[theorem]{Proposition}
\newtheorem{lemma}[theorem]{Lemma}
\newtheorem{corollary}[theorem]{Corollary}
\theoremstyle{definition}
\newtheorem{definition}[theorem]{Definition}
\newtheorem{example}[theorem]{Example}
\theoremstyle{remark}
\newtheorem{remark}[theorem]{Remark}
\newcommand{\R}{\mathbb R}
\newcommand{\Q}{\mathbb Q}
\newcommand{\N}{\mathbb Z_{\geq0}}
\newcommand{\tHilb}{\operatorname{TH}}
\newcommand{\degT}{\operatorname{deg}_{T}}
\newcommand{\ceil}[1]{\left\lceil #1\right\rceil}
\newcommand{\floor}[1]{\left\lfloor #1\right\rfloor}
\newcommand{\ip}[2]{\langle #1,#2\rangle}
\newcommand{\VD}{V_d}
\newcommand{\HD}{H_d}
\newcommand{\BoxH}{\tHilb^{\Box}}

\title{\textbf{Nonintegral Tropical Degree\\from Two Horizontal Segments}\\[0.45em]
\large An exact Hilbert function, explicit rational certificates,\\and counterexamples to eventual polynomiality}
\author{AI-assisted research draft}
\date{20 September 2026}

\begin{document}
\maketitle
\begin{abstract}
For Grigoriev's strict-witness tropical Hilbert function, with monomials of
\emph{total degree at most} $k$, we compute the entire function of a compact
one-dimensional min-plus prevariety consisting of two horizontal segments.
For every positive rational $d$, let
\[
V_d=([0,1]\times\{0\})\ \cup\ ([d+1,d+2]\times\{1\}).
\]
We prove, for every integer $k\geq0$,
\[
\tHilb_{V_d}(k)=k+\min\{k,\lceil k/d\rceil\}+1.
\]
In particular, $V_3$ has tropical degree $4/3$, and its Hilbert function is a
quasipolynomial of least period three, not an eventual polynomial. These give
negative answers, under the definitions in arXiv:2404.06440v2, to the
integrality question and the eventual-polynomiality conjecture stated in that
paper's introduction. Every rational number strictly between one and two is
realized as a tropical degree by the same family. For $d\geq1$, the formula
also holds with a box degree cutoff. The proof uses only ordered slopes,
a strict inequality between facing endpoints, and an explicit family of
quadratic coefficient profiles. All lower bounds have rational strict
witnesses. An accompanying standard-library Python verifier checks 823,974
witness inequalities in exact arithmetic, in addition to endpoint bounds and
recurrence identities.
\end{abstract}

\medskip
\noindent\textbf{Status.} The mathematical claims below have complete proofs in
this draft. The proposed resolution has not been independently refereed or
checked in a proof assistant. The source and novelty audit records what was
actually retrieved on 20 September 2026; it is not an exhaustive priority
certification.

\clearpage
\begingroup
\small
\setlength{\parskip}{0pt}
\tableofcontents
\endgroup
\clearpage

\section{The problem and the result}
\label{sec:problem}

\subsection{The precise published target}

The target is the tropical Hilbert function in Definitions 1.1 and 1.2 of
Dima Grigoriev's \emph{Degree of 1-dimensional tropical prevariety},
arXiv:2404.06440v2, revised 11 August 2025 \cite{Grigoriev2025}.
This function counts how many bounded-degree monomials can each be the unique
minimum somewhere on a fixed set, using \emph{one common assignment of
coefficients}. Its definition is reproduced carefully in
Section~\ref{sec:definitions}.

In the introduction of that version, the author conjectures eventual
polynomiality of this function and asks, as Question 2, whether the tropical
degree is always an integer. The paper distinguishes total-degree and
box-degree cutoffs. Our main result concerns total degree; we separately
establish the relevant box-degree result in Section~\ref{sec:box}.
The version matters: the original 2024 submission had the different title
\emph{A tropical version of Hilbert function} \cite{Grigoriev2024}.

\subsection{An all-degree formula}

For $d\in\Q_{>0}$ put
\begin{equation}
\label{eq:V}
I_0=[0,1]\times\{0\},\qquad
I_1=[d+1,d+2]\times\{1\},\qquad
V_d=I_0\cup I_1.
\end{equation}
Both components have length one. Their heights differ by one, and the
horizontal displacement from the right endpoint of $I_0$ to the left
endpoint of $I_1$ is $d$.

\begin{theorem}[Exact Hilbert function]
\label{thm:main}
For every $d\in\Q_{>0}$ and every $k\in\N$,
\begin{equation}
\label{eq:main}
H_d(k):=\tHilb_{V_d}(k)
=k+\min\left\{k,\ceil{\frac{k}{d}}\right\}+1.
\end{equation}
In particular,
\begin{equation}
\label{eq:piecewise}
H_d(k)=
\begin{cases}
2k+1,&0<d\leq1,\\[2pt]
k+\lceil k/d\rceil+1,&d\geq1.
\end{cases}
\end{equation}
The two expressions agree when $d=1$.
\end{theorem}

\begin{corollary}[A counterexample with tropical degree $4/3$]
\label{cor:main}
For
\[
V_3=([0,1]\times\{0\})\cup([4,5]\times\{1\}),
\]
one has
\begin{equation}
\label{eq:d3}
\tHilb_{V_3}(k)=k+\ceil{k/3}+1,
\qquad
\degT(V_3)=\lim_{k\to\infty}\frac{\tHilb_{V_3}(k)}{k}=\frac43.
\end{equation}
The Hilbert function is not eventually a polynomial.
\end{corollary}

The full proof of Theorem~\ref{thm:main} is given in
Sections~\ref{sec:upper} and~\ref{sec:lower}. No asymptotic existence theorem,
tropical rank theorem, or general classification of prevarieties is required
for that proof. We give min-plus equations for the set itself, so membership
in the source's category is explicit rather than assumed.

\begin{figure}[ht]
\centering
\begin{tikzpicture}[x=1.7cm,y=1.45cm,>=Stealth]
\draw[->,thin] (-.25,-.35)--(5.45,-.35) node[right] {$x$};
\draw[->,thin] (-.25,-.35)--(-.25,1.4) node[above] {$y$};
\draw[line width=1.6pt] (0,0)--(1,0);
\draw[line width=1.6pt] (4,1)--(5,1);
\foreach \x/\y in {0/0,1/0,4/1,5/1}{\fill (\x,\y) circle(1.5pt);}
\node[above] at (.5,0) {$I_0$};
\node[above] at (4.5,1) {$I_1$};
\node[below] at (0,0) {$0$};\node[below] at (1,0) {$1$};
\node[below] at (4,1) {$(4,1)$};\node[below] at (5,1) {$(5,1)$};
\draw[dashed,thin,->] (1,.07)--(3.94,.98);
\node[above left] at (2.55,.55) {$(d,1)=(3,1)$};
\end{tikzpicture}
\caption{The counterexample $V_3$. The dashed arrow is an endpoint displacement,
not an additional component of the prevariety.}
\label{fig:geometry}
\end{figure}

\FloatBarrier
\section{Definitions and conventions}
\label{sec:definitions}

\subsection{Monomials, coefficients, and the total-degree cutoff}

In min-plus notation, an ordinary affine function
\[
c+ix+jy,\qquad i,j\in\N,\quad c\in\R,
\]
is a tropical monomial with coefficient $c$ and total degree $i+j$.
A tropical polynomial is the minimum of finitely many such affine functions.
We use the finite exponent set
\begin{equation}
\label{eq:Delta}
\Delta_k=\{(i,j)\in\N^2:i+j\leq k\}.
\end{equation}
The coefficient-free monomial corresponding to $u=(i,j)$ is the function
$x,y\mapsto ix+jy$. Additive constants are chosen when constructing a
certificate; they do not create new exponent vectors.

\begin{definition}[Strict-witness independence]
\label{def:independent}
A finite set $S\subseteq\Delta_k$ is independent on $V\subseteq\R^2$ if there
exist coefficients $(c_u)_{u\in S}\in\R^S$ and points $(z_u)_{u\in S}\in V^S$
such that
\begin{equation}
\label{eq:witness}
c_u+\ip{u}{z_u}<c_v+\ip{v}{z_u}
\quad\text{for every }u\in S\text{ and }v\in S\setminus\{u\}.
\end{equation}
The tropical Hilbert function is
\begin{equation}
\label{eq:THdefinition}
\tHilb_V(k)=\max\{|S|:S\subseteq\Delta_k\text{ is independent on }V\}.
\end{equation}
\end{definition}

This is exactly the strict-inequality notion in the cited Definitions 1.1
and 1.2 \cite{Grigoriev2025}. We do not replace it by another notion of tropical
linear independence. In particular, the coefficient vector in
\eqref{eq:witness} is fixed while the witness point varies. Choosing a separate
coefficient vector at each point would define a different and much weaker
condition.

The maximum in \eqref{eq:THdefinition} is over subsets of a finite set, so it
is a genuine maximum. On any nonempty $V$, one has $\tHilb_V(0)=1$: the sole
exponent vector is $(0,0)$, whose singleton support is independent.

For a one-dimensional set for which the limit exists, the degree considered
here is
\begin{equation}
\label{eq:degree}
\degT(V)=\lim_{k\to\infty}\frac{\tHilb_V(k)}{k}.
\end{equation}
Our exact formula will establish this limit directly for every $V_d$.

\subsection{What the result does not identify}

An ideal in the sense of Maclagan--Rinc\'on has additional tropical linear-space
structure \cite{MaclaganRincon}. Its ideal-theoretic Hilbert function is not
being substituted for Definition~\ref{def:independent}. Likewise, a compact
union of two isolated segments is not asserted to be a balanced tropical
curve. The balancedness theorem for varieties of tropical ideals belongs to
a different, more restrictive setting \cite{Balanced}.

The counterexample addresses the class actually specified in
\cite{Grigoriev2025}: min-plus prevarieties, including finite unions of rational
polyhedra. It makes no claim against an integrality theorem whose hypotheses
include balancing, an ideal structure, or a different normalization.

\section{Explicit equations for the prevariety}
\label{sec:equations}

\begin{proposition}[A three-equation presentation]
\label{prop:equations}
Write $d=p/q$ with positive integers $p,q$. The set $V_d$ is precisely the
solution set in $\R^2$ of
\begin{align}
\min\{2y,1\}&=y,\label{eq:eq1}\\
\min\{qx,(p+q)y\}&=(p+q)y,\label{eq:eq2}\\
\min\{qx,(p+q)y+q\}&=qx.\label{eq:eq3}
\end{align}
Every expression on either side is a tropical polynomial with nonnegative
integer exponent vectors.
\end{proposition}

\begin{proof}
If $2y\leq1$, equation~\eqref{eq:eq1} forces $2y=y$, and hence $y=0$.
If $2y\geq1$, it forces $1=y$, and hence $y=1$. Conversely, both $y=0$ and
$y=1$ satisfy it. Equations~\eqref{eq:eq2} and~\eqref{eq:eq3} are respectively
equivalent to
\[
qx\geq(p+q)y,\qquad qx\leq(p+q)y+q.
\]
Thus $y=0$ gives $0\leq x\leq1$, whereas $y=1$ gives
$d+1\leq x\leq d+2$. These are exactly the two components in \eqref{eq:V}.
All coefficients multiplying $x$ and $y$ in the displayed tropical monomials
are nonnegative integers; their additive constants are rational.
\end{proof}

For the main counterexample $p=3,q=1$, the presentation is especially small:
\begin{equation}
\label{eq:eqd3}
\min\{2y,1\}=y,\qquad
\min\{x,4y\}=4y,\qquad
\min\{x,4y+1\}=x.
\end{equation}
The set is nonempty, compact, one-dimensional, and has two connected
components. These are direct properties of the displayed segments.

\subsection{Why a union creates a new obstruction}

On a single horizontal interval, the restriction of an exponent $(i,j)$ has
slope $i$. Two distinct monomials with the same slope cannot both be unique
minima on that interval. There are only $k+1$ possible slopes, and it is easy
to realize all of them.

Two intervals give two opportunities to use the same slope with different
vertical exponents. Those opportunities are not independent: the same
coefficient vector must work on both components. Moving between the facing
endpoints imposes a strict comparison between two monomials. The horizontal
gap $d$ then converts the finite vertical-exponent budget into a bound on the
number of additional slopes that can occur. This bound contains a ceiling,
and the ceiling persists at all degrees.

\section{The upper bound}
\label{sec:upper}

We first isolate the elementary facts that make the proof insensitive to
degeneracies and ties.

\subsection{Perturbing coefficients without losing witnesses}

\begin{lemma}[Generic endpoint winners]
\label{lem:perturb}
Suppose $S$ is independent on $V_d$. Its coefficients can be chosen so that
all members of $S$ still have strict witnesses and the minimum is attained by
a unique member of $S$ at each of the four endpoints of the two segments.
\end{lemma}

\begin{proof}
For a fixed certificate, every difference in \eqref{eq:witness} is positive.
When $|S|\geq2$, let $\gamma$ be the minimum of these finitely many positive
differences. A perturbation of every coefficient by less than $\gamma/4$
changes any difference by less than $\gamma/2$, so it preserves every witness.

At a fixed endpoint $z$, a tie between distinct $u,v\in S$ is the proper affine
hyperplane
\[
c_u-c_v=\ip{v-u}{z}
\]
in coefficient space. The four endpoints give finitely many such hyperplanes.
Their union contains no nonempty open ball. We can therefore perturb within
the permitted ball while avoiding all endpoint ties. A singleton support
requires no perturbation.
\end{proof}

A monomial with a strict witness at an endpoint is also uniquely minimal on
a sufficiently short one-sided interval, by continuity. Thus endpoint
witnesses do not introduce exceptional zero-length pieces into the slope
count below.

\subsection{Counting strictly active slopes}

Call a monomial \emph{active on a segment} if it is uniquely minimal at some
point of that segment. Under the generic endpoint convention, the set of
active monomials includes the endpoint winners.

\begin{lemma}[Slope order on a horizontal segment]
\label{lem:slopes}
Along a horizontal segment traversed from left to right, the $x$-slopes of
successive distinct active monomials strictly decrease. In particular, if the
left and right endpoint slopes are $u$ and $v$, the number of active monomials
is at most $u-v+1$.
\end{lemma}

\begin{proof}
Let two distinct monomials restrict to $\alpha+it$ and $\beta+jt$.
Suppose the first wins strictly at $t_1$ and the second at $t_2>t_1$. Then
\[
\alpha+it_1<\beta+jt_1,
\qquad
\beta+jt_2<\alpha+it_2.
\]
Adding gives $(i-j)(t_2-t_1)>0$, so $i>j$. Distinct monomials of the same
slope cannot both be strictly active: their restricted difference is
constant. Every active integer slope lies between the endpoint slopes, and
there are $u-v+1$ integers in that interval.
\end{proof}

Now fix an independent support $S\subseteq\Delta_k$ and coefficients supplied
by Lemma~\ref{lem:perturb}. Let $S_0,S_1$ be the sets active on $I_0,I_1$,
respectively. Because every element of $S$ has a witness,
\begin{equation}
\label{eq:cover}
S=S_0\cup S_1.
\end{equation}
Write the four endpoint slopes in order as
\[
\begin{array}{c|cc}
&\text{left endpoint}&\text{right endpoint}\\ \hline
I_0&u&r\\
I_1&P&v.
\end{array}
\]
Lemma~\ref{lem:slopes} gives
\begin{equation}
\label{eq:counts}
|S_0|\leq u-r+1,\qquad |S_1|\leq P-v+1,
\qquad 0\leq v,r,P,u\leq k.
\end{equation}

\begin{lemma}[The universal two-segment bound]
\label{lem:2k}
For total degree at most $k$,
\[
|S|\leq2k+1.
\]
\end{lemma}

\begin{proof}
Each active set has size at most $k+1$. If either has size at most $k$, the
claim follows immediately from \eqref{eq:cover}. Otherwise both have size
$k+1$, so each contains a monomial of $x$-slope $k$. The only exponent in
$\Delta_k$ with that slope is $(k,0)$. Consequently the two sets have a common
element, and their union has size at most $2(k+1)-1$.
\end{proof}

This is the one place in the upper-bound proof where the triangular degree
cutoff does more than bound each exponent separately. With a box cutoff,
slope $k$ no longer identifies a unique exponent vector.

\subsection{Comparing the facing endpoints}

Let $a=(r,s)$ be the exponent winning at $z=(1,0)$, and let $b=(P,Q)$ be the
exponent winning at $w=(d+1,1)$. If $a=b$, then $P=r$ and $S_0\cap S_1$ is
nonempty. From \eqref{eq:counts},
\begin{equation}
\label{eq:shared}
|S|\leq(u-r+1)+(r-v+1)-1=u-v+1\leq k+1.
\end{equation}
This covers the shared-winner case completely.

Suppose next that $a\ne b$. Strictness at the two endpoints gives
\[
c_a+\ip{a}{z}<c_b+\ip{b}{z},
\qquad
c_b+\ip{b}{w}<c_a+\ip{a}{w}.
\]
Adding cancels the coefficients and yields
\begin{equation}
\label{eq:facing}
0<\ip{a-b}{w-z}=d(r-P)+s-Q.
\end{equation}
Since $s\leq k-r$ and $Q\geq0$,
\begin{equation}
\label{eq:budget}
d(P-r)<s-Q\leq k-r\leq k.
\end{equation}
The integer $P-r$ is therefore strictly less than $k/d$, so
\begin{equation}
\label{eq:ceiling}
P-r\leq\ceil{k/d}-1.
\end{equation}
This statement is valid even if $P-r$ is negative. Using
\eqref{eq:counts} and $u\leq k$, $v\geq0$, we obtain
\begin{align}
|S|&\leq |S_0|+|S_1|\nonumber\\
&\leq u-r+1+P-v+1\nonumber\\
&\leq k+(P-r)+2\nonumber\\
&\leq k+\ceil{k/d}+1.\label{eq:gapupper}
\end{align}

\begin{proposition}[Upper half of the main theorem]
\label{prop:upper}
For all $k\geq0$,
\[
\tHilb_{V_d}(k)\leq k+\min\{k,\lceil k/d\rceil\}+1.
\]
\end{proposition}

\begin{proof}
For $k=0$ the assertion was already noted. For $k\geq1$, the shared-winner
bound \eqref{eq:shared} is no larger than the claimed upper bound. In the
other case combine \eqref{eq:gapupper} with Lemma~\ref{lem:2k}.
\end{proof}

\begin{remark}[Why the strict ceiling is essential]
The implication $n<t\Rightarrow n\leq\lceil t\rceil-1$ is the source of the
periodic term. Replacing \eqref{eq:facing} by a non-strict inequality would lose
one unit whenever $k/d$ is an integer. The generic perturbation does not create
this strictness artificially: it preserves an already strict certificate
while ensuring that the selected endpoints also have unique winners.
\end{remark}

\section{An explicit lower-bound certificate}
\label{sec:lower}

We now construct a support of exactly the size in Proposition~\ref{prop:upper}.
The construction gives every coefficient and every witness as a rational
number. It is therefore stronger than an existence argument based on a
linear-programming feasibility test.

\subsection{The parameters}

Fix $k\geq1$ and $d\in\Q_{>0}$. Define
\begin{equation}
\label{eq:mwsigma}
m=\min\{k,\lceil k/d\rceil\}-1,\qquad
\omega=\min\{1,d\},\qquad
\sigma=\omega k-dm.
\end{equation}
Then $0\leq m<k$ and $\sigma>0$. To check the last assertion, if $d\geq1$,
then $m=\lceil k/d\rceil-1<k/d$, and hence $\sigma=k-dm>0$. If $d\leq1$,
then $m=k-1$, so $\sigma=dk-d(k-1)=d>0$.

Choose the small positive rational numbers
\begin{equation}
\label{eq:small}
\delta=\min\left\{\frac18,\frac{\sigma}{8(m+1)}\right\},\qquad
\varepsilon=\frac{\delta}{4k},\qquad
\eta=\frac{\delta}{4(m+1)},
\end{equation}
and set
\begin{equation}
\label{eq:abh}
a=1-\delta,\qquad b=d+1+\delta,\qquad h=\omega k-\delta.
\end{equation}
The letters $a,b$ in this section are scalar parameters, not the endpoint
exponent vectors used in the upper bound.

\subsection{Two groups of monomials}

Take the following coefficient-bearing monomials:
\begin{align}
A_i(x,y)&=\varepsilon i^2-ai+ix+(k-i)y,
&&0\leq i\leq k,\label{eq:Ai}\\
B_j(x,y)&=h-bj+\eta j^2+jx,
&&0\leq j\leq m.\label{eq:Bj}
\end{align}
Thus the exponent sets are the diagonal row
\[
\{(i,k-i):0\leq i\leq k\}
\]
and a short initial part of the bottom row
\[
\{(j,0):0\leq j\leq m\}.
\]
They are disjoint because $m<k$: the sole diagonal exponent with second
coordinate zero is $(k,0)$, which is not in the bottom-row group.
Every exponent lies in $\Delta_k$.

The total support size is
\begin{equation}
\label{eq:size}
(k+1)+(m+1)=k+\min\{k,\lceil k/d\rceil\}+1.
\end{equation}
We must show that the minimum of all the functions in
\eqref{eq:Ai}--\eqref{eq:Bj} is a certificate of independence.

\subsection{Witnesses for the diagonal row}

For $0\leq i\leq k$, use
\begin{equation}
\label{eq:Awit}
z_i=(a-2\varepsilon i,0).
\end{equation}
Its horizontal coordinate satisfies
\[
1-\frac32\delta\leq a-2\varepsilon i\leq1-\delta.
\]
Because $0<\delta\leq1/8$, each $z_i$ is strictly inside $I_0$.
For another diagonal-row index $\ell$,
\begin{align}
A_\ell(z_i)
&=\varepsilon\ell^2-a\ell+\ell(a-2\varepsilon i)\nonumber\\
&=\varepsilon(\ell-i)^2-\varepsilon i^2.\label{eq:Asquare}
\end{align}
Consequently $A_i(z_i)=-\varepsilon i^2\leq0$, and $A_i$ is the unique
minimum within the $A$ group.

It remains to keep all $B$ terms above it. Put
$x_{\min}=1-3\delta/2$. For $0\leq j\leq m$, using the nonnegativity of
$\eta j^2$ gives
\begin{align}
B_j(z_i)
&=h-(b-z_{i,x})j+\eta j^2\nonumber\\
&\geq h-(b-x_{\min})m\nonumber\\
&=\omega k-\delta-(d+5\delta/2)m\nonumber\\
&=\sigma-\delta(1+5m/2).\label{eq:Bpositive}
\end{align}
Here $b-z_{i,x}>0$, so replacing $j$ by $m$ in the negative term is legitimate.
The choice of $\delta$ implies
\[
\delta(1+5m/2)
\leq \sigma\frac{1+5m/2}{8(m+1)}
\leq\frac{5\sigma}{16}.
\]
Therefore
\begin{equation}
\label{eq:Bmargin}
B_j(z_i)\geq\frac{11\sigma}{16}>0\geq A_i(z_i).
\end{equation}
Every $A_i$ has a strict witness against \emph{all} other monomials.

\subsection{Witnesses for the bottom row}

For $0\leq j\leq m$, use
\begin{equation}
\label{eq:Bwit}
w_j=(b-2\eta j,1).
\end{equation}
Their horizontal coordinates satisfy
\[
d+1+\frac\delta2<b-2\eta j\leq d+1+\delta<d+2.
\]
(The weaker left inequality with $\leq$ would also suffice.) Thus the witnesses
are strictly inside $I_1$.
For another bottom-row index $\ell$,
\begin{equation}
\label{eq:Bsquare}
B_\ell(w_j)=h+\eta(\ell-j)^2-\eta j^2.
\end{equation}
It follows that $B_j$ is the unique minimum within its own group, and
$B_j(w_j)=h-\eta j^2\leq h$.

For any $0\leq i\leq k$,
\begin{equation}
\label{eq:Alower}
A_i(w_j)=k+i(w_{j,x}-1-a)+\varepsilon i^2.
\end{equation}
Since $w_{j,x}-1-a\geq d-1+3\delta/2$, if $d\geq1$ then
$A_i(w_j)\geq k=\omega k$. If $0<d\leq1$, then
\[
A_i(w_j)\geq k+i(d-1)\geq k+k(d-1)=dk=\omega k,
\]
where the second inequality uses $i\leq k$ and $d-1\leq0$.
In both cases,
\begin{equation}
\label{eq:Amargin}
A_i(w_j)\geq\omega k=h+\delta>B_j(w_j).
\end{equation}
Thus every $B_j$ also has a strict witness against every competitor.

\begin{proposition}[Lower half of the main theorem]
\label{prop:lower}
The tropical polynomial
\[
F_{k,d}(x,y)=\min\{A_0(x,y),\ldots,A_k(x,y),B_0(x,y),\ldots,B_m(x,y)\}
\]
is a certificate of independence for its full support. Hence
\[
\tHilb_{V_d}(k)\geq k+\min\{k,\lceil k/d\rceil\}+1.
\]
\end{proposition}

\begin{proof}
Equations~\eqref{eq:Asquare}, \eqref{eq:Bmargin}, \eqref{eq:Bsquare}, and
\eqref{eq:Amargin} verify all the strict inequalities in
Definition~\ref{def:independent}. The support count is \eqref{eq:size}.
\end{proof}

\begin{proof}[Proof of Theorem~\ref{thm:main}]
For $k\geq1$, combine Propositions~\ref{prop:upper} and~\ref{prop:lower}.
For $k=0$, the unique coefficient-free monomial has a singleton certificate,
so the formula also holds.
\end{proof}

\subsection{Quantitative strictness and certificate size}

The construction supplies useful positive margins. Within the $A$ group,
the gap at a witness is at least $\varepsilon$; within the $B$ group, when it
has more than one member, the gap is at least $\eta$. Across the groups,
\eqref{eq:Bmargin} gives at least $11\sigma/16$ for an $A$ witness, and
\eqref{eq:Amargin} gives at least $\delta$ for a $B$ witness. Thus, for $k\geq1$,
a valid uniform lower bound for all competitor gaps is
\begin{equation}
\label{eq:margin}
\min\{\varepsilon,\eta,\delta,11\sigma/16\}>0.
\end{equation}
Including $\eta$ in this minimum is harmless when the $B$ group is a singleton.

There are at most $2k+1$ monomials, so direct certificate verification needs
at most $(2k+1)(2k)$ ordered strict comparisons. The generator uses $O(k)$
rational arithmetic operations. For $d=p/q$ in lowest terms, the displayed
formulas give rational numbers with bit lengths bounded by
$O(\log(k+1)+\log(p+1)+\log(q+1))$ per coefficient or coordinate. Indeed,
$m\leq k$, and the denominators are products of $q$ and a bounded number of
factors among $8$, $4k$, and $m+1$; the numerators grow at most polynomially
in $p,q,k$. This is an explicit, polynomial-size certificate, not a limiting
or numerical construction.

\section{Degree, quasipolynomiality, and generating functions}
\label{sec:consequences}

\subsection{Varying the gap and the degree}

\begin{corollary}[Degree formula]
\label{cor:degree}
For every positive rational $d$,
\begin{equation}
\label{eq:degreeformula}
\degT(V_d)=1+\min\{1,1/d\}.
\end{equation}
Every rational $\rho$ with $1<\rho\leq2$ is realized by a member of this
family. The denominators of these degrees are unbounded, even though the
number, lengths, and directions of the components are fixed.
\end{corollary}

\begin{proof}
Divide \eqref{eq:piecewise} by $k$ and use
$0\leq\lceil k/d\rceil-k/d<1$. For $1<\rho<2$, take
$d=1/(\rho-1)>1$. For $\rho=2$, take $d=1$.
For example, $d=N$ gives degree $(N+1)/N$ for every integer $N\geq2$, and
these fractions are in lowest terms.
\end{proof}

The value $1$ is approached as $d\to\infty$, but is not attained for finite
$d$ in this two-segment family. A single nondegenerate horizontal segment
has Hilbert function $k+1$, and hence degree one: the upper slope count is
sharp by the same quadratic-profile construction used above. Thus the
endpoint value itself is not problematic; it simply has a different witness.

\subsection{Exact quasipolynomial periods}

\begin{definition}
A function $f:\N\to\R$ is an eventual quasipolynomial of period $t\geq1$ if,
for each residue class modulo $t$, it agrees with a polynomial in $k$ for all
sufficiently large $k$ in that class. An eventual polynomial is the special
case $t=1$.
\end{definition}

\begin{corollary}[Least period]
\label{cor:period}
Let $d=p/q>1$ in lowest terms, so $p>q\geq1$. Then
\begin{equation}
\label{eq:rationalH}
H_d(k)=k+\ceil{\frac{qk}{p}}+1,
\qquad
H_d(k+p)=H_d(k)+p+q.
\end{equation}
It is a quasipolynomial of least eventual period $p$. In particular,
it is not an eventual polynomial. Within the rational family $V_d$, eventual
polynomiality holds exactly when $0<d\leq1$.
\end{corollary}

\begin{proof}
Write $k=pn+r$, $0\leq r<p$. Formula~\eqref{eq:rationalH} becomes
\begin{equation}
\label{eq:residuegeneral}
H_d(pn+r)=(p+q)n+r+\ceil{\frac{qr}{p}}+1.
\end{equation}
This proves quasipolynomiality with period $p$ and the shift identity.

To prove minimality, set
\[
E(k)=H_d(k)-\left(1+\frac qp\right)k-1
=\ceil{\frac{qk}{p}}-\frac{qk}{p}.
\]
This bounded function is nonnegative and vanishes exactly when $p$ divides
$k$. Suppose $H_d$ has an eventual quasipolynomial period $t$. On each residue
class modulo $t$, its polynomial must be linear with slope $1+q/p$:
boundedness of $E(k)$ excludes both a higher degree and any different slope.
Therefore $E(k)$ is eventually periodic with period $t$. Choose a sufficiently
large multiple $k$ of $p$. Since $E(k)=0$, eventual periodicity gives
$E(k+t)=0$, which implies $p\mid t$. Thus the least eventual period is $p$.
The assertion for $d\leq1$ follows from $H_d(k)=2k+1$.
\end{proof}

A shorter argument suffices merely to disprove eventual polynomiality.
An eventual polynomial with linear growth must be affine. If it has integer
values at all sufficiently large integers, its slope, being a difference
of consecutive integer values, is an integer. Here the slope $1+q/p$ lies
strictly between one and two, so that is impossible.

\subsection{The complete sequence for \texorpdfstring{$d=3$}{d = 3}}

\begin{corollary}
\label{cor:d3sequence}
Writing $k=3n+r$, the main example has
\begin{align*}
H_3(3n)&=4n+1,\\
H_3(3n+1)&=4n+3,\\
H_3(3n+2)&=4n+4.
\end{align*}
Its first values are
\[
1,3,4,5,7,8,9,11,12,13,15,16,17,19,20,21,\ldots.
\]
The first differences repeat $2,1,1$ from the start.
\end{corollary}

\begin{proof}
Substitute the three residue classes into \eqref{eq:d3}.
\end{proof}

The ordinary generating function has a short reduced rational form:
\begin{equation}
\label{eq:GF}
G_3(z)=\sum_{k\geq0}H_3(k)z^k
=\frac{(1+z)^2}{(1-z)(1-z^3)}.
\end{equation}
For completeness, put $C(k)=\lceil k/3\rceil$. Then $C(0)=0$ and
$C(k)-C(k-1)$ equals one precisely when $k\equiv1\pmod3$. Consequently
\[
(1-z)\sum_{k\geq0}C(k)z^k=\frac{z}{1-z^3},
\]
and hence
\[
G_3(z)=\frac{1}{(1-z)^2}+\frac{z}{(1-z)(1-z^3)}
=\frac{1+2z+z^2}{(1-z)(1-z^3)}.
\]
The numerator is nonzero at every cube root of unity, including one, so the
displayed fraction is reduced. Multiplying by the denominator gives the
constant-coefficient recurrence
\begin{equation}
\label{eq:recurrence}
H_3(k)=H_3(k-1)+H_3(k-3)-H_3(k-4),\qquad k\geq4,
\end{equation}
with initial values $1,3,4,5$. Rationality of a generating function is thus
fully compatible with failure of eventual polynomiality of its coefficients.

\subsection{A rational generating function for every rational gap}

For $d=p/q>1$ in lowest terms, define the integers
\[
e_r=\ceil{\frac{qr}{p}}-\ceil{\frac{q(r-1)}{p}},\qquad 1\leq r\leq p.
\]
They belong to $\{0,1\}$ and sum to $q$. The increments of $\lceil qk/p\rceil$
repeat this block. Applying the same first-difference argument gives
\begin{equation}
\label{eq:GFgeneral}
\sum_{k\geq0}H_d(k)z^k=
\frac{1+z+\cdots+z^{p-1}+\sum_{r=1}^{p}e_rz^r}
{(1-z)(1-z^p)}.
\end{equation}
This formula is not asserted to be reduced for every $p,q$. Its role is to
give a direct rational representation and another way to generate the
sequence exactly. When $0<d\leq1$, the generating function is simply
$(1+z)/(1-z)^2$.

\section{The separate box-degree convention}
\label{sec:box}

The source also considers
\[
\Box_k=\{(i,j)\in\N^2:0\leq i\leq k,\ 0\leq j\leq k\}
\]
in place of $\Delta_k$ \cite{Grigoriev2025}. Let $\BoxH_V(k)$ denote the resulting
maximum number of independently witnessed monomials. It is important to
check this modification rather than silently identifying it with total degree.

\begin{proposition}[The same counterexample for a box cutoff]
\label{prop:box}
For every rational $d\geq1$ and $k\geq0$,
\begin{equation}
\label{eq:box}
\BoxH_{V_d}(k)=\tHilb_{V_d}(k)=k+\lceil k/d\rceil+1.
\end{equation}
Thus the box-normalized degree of $V_3$ is also $4/3$, and its box Hilbert
function is not eventually polynomial.
\end{proposition}

\begin{proof}
The inclusion $\Delta_k\subseteq\Box_k$ and the existing certificate give the
lower bound. For the upper bound, apply the perturbation and slope lemmas to
an independent subset of $\Box_k$. The shared-winner case still gives
$|S|\leq k+1$.

For distinct facing winners, \eqref{eq:facing} still holds. Although we no
longer have $s\leq k-r$, we still have $s\leq k$ and $Q\geq0$. Therefore
$d(P-r)<k$, which is all that is needed for \eqref{eq:ceiling} and
\eqref{eq:gapupper}. It follows that $|S|\leq k+\lceil k/d\rceil+1$.
For $d\geq1$, this agrees with the total-degree lower bound. The case $k=0$
again consists of the unique constant monomial.
\end{proof}

No box formula for $d<1$ is needed for either counterexample, and none is
being inferred from the total-degree formula in that range. In particular,
Lemma~\ref{lem:2k} itself is not a valid box-degree argument.

\section{A seven-monomial certificate that can be checked by hand}
\label{sec:example}

The first few values already expose the nonconstant increments, but a small
explicit certificate makes the lower-bound mechanism tangible. Take $d=3$
and $k=4$. The theorem predicts $H_3(4)=7$.

Use five diagonal-row monomials
\[
A_i(x,y)=\frac{i^2}{100}-\frac{9i}{10}+ix+(4-i)y,
\qquad 0\leq i\leq4,
\]
together with
\[
B_0(x,y)=\frac{39}{10},\qquad B_1(x,y)=-\frac{19}{100}+x.
\]
These parameters are chosen for a compact table; they are not the particular
small-parameter choice in \eqref{eq:small}. Their strict witnesses are listed
below.

\begin{table}[ht]
\centering
\renewcommand{\arraystretch}{1.22}
\begin{tabular}{@{}c c r c r@{}}
\toprule
Term & Exponent & Coefficient & Witness & Own value\\
\midrule
$A_0$ & $(0,4)$ & $0$ & $(9/10,0)$ & $0$\\
$A_1$ & $(1,3)$ & $-89/100$ & $(22/25,0)$ & $-1/100$\\
$A_2$ & $(2,2)$ & $-44/25$ & $(43/50,0)$ & $-4/100$\\
$A_3$ & $(3,1)$ & $-261/100$ & $(21/25,0)$ & $-9/100$\\
$A_4$ & $(4,0)$ & $-86/25$ & $(41/50,0)$ & $-16/100$\\
$B_0$ & $(0,0)$ & $39/10$ & $(41/10,1)$ & $390/100$\\
$B_1$ & $(1,0)$ & $-19/100$ & $(102/25,1)$ & $389/100$\\
\bottomrule
\end{tabular}
\caption{An exact seven-term certificate at $k=4$. Every witness belongs to
$V_3$, and every exponent has total degree at most four.}
\label{tab:certificate}
\end{table}

Let $T_u$ denote a row's selected term and $z_u$ its witness. To avoid relying
on a plotted lower envelope or rounded values, Table~\ref{tab:gaps} prints
\emph{all} competitor gaps, multiplied by $100$:
\[
G_{u,v}=100\bigl(T_v(z_u)-T_u(z_u)\bigr).
\]
The diagonal is zero and every off-diagonal entry is a positive integer.
This proves all 42 strict comparisons directly, with least gap $1/100$.

\begin{table}[ht]
\centering
\renewcommand{\arraystretch}{1.15}
\begin{tabular}{@{}c rrrrrrr@{}}
\toprule
Witness for & $A_0$ & $A_1$ & $A_2$ & $A_3$ & $A_4$ & $B_0$ & $B_1$\\
\midrule
$A_0$ &0&1&4&9&16&390&71\\
$A_1$ &1&0&1&4&9&391&70\\
$A_2$ &4&1&0&1&4&394&71\\
$A_3$ &9&4&1&0&1&399&74\\
$A_4$ &16&9&4&1&0&406&79\\
$B_0$ &10&231&454&679&906&0&1\\
$B_1$ &11&230&451&674&899&1&0\\
\bottomrule
\end{tabular}
\caption{The complete scaled gap matrix for Table~\ref{tab:certificate}.
Columns identify competitors, not alternative witnesses.}
\label{tab:gaps}
\end{table}

The endpoint upper bound gives $H_3(4)\leq4+\lceil4/3\rceil+1=7$.
The small example is therefore optimal, not merely a nonmaximal independent
set. Its JSON representation is included in the accompanying data.

\section{Exact computation and reproducibility}
\label{sec:computation}

\subsection{What was actually checked}

The Python program \texttt{code/verify.py} uses only the standard library,
and all numerical values relevant to certification are instances of
\texttt{fractions.Fraction}. No floating-point tolerances, numerical linear
programming, or external computer algebra system enter the verifier.

The recorded run used every integer $0\leq k\leq40$ and the fourteen values
\[
\frac14,\ \frac12,\ \frac23,\ 1,\ \frac54,\ \frac32,\ \frac53,\ 2,\
\frac73,\ 3,\ 4,\ 5,\ 10,\ \frac{101}{7}.
\]
For each case it constructed the exact certificate and checked that all
exponents were distinct, their degrees satisfied the cutoff, all witnesses
lay in the correct segments, and every selected term was strictly smaller
than every other selected term at its witness.

\begin{table}[ht]
\centering
\begin{tabular}{@{}lr@{}}
\toprule
Recorded check & Count or result\\
\midrule
Rational gap parameters & 14\\
Parameter--degree cases & 574\\
Strict witness inequalities, including the hand example & 823,974\\
Total-degree endpoint-relaxation checks & 574\\
Box endpoint-relaxation checks for $d\geq1$ & 451\\
Generating-function coefficients checked for $d=3$ & 121\\
Least competitor gap in the hand example & $1/100$\\
Overall result & PASS\\
\bottomrule
\end{tabular}
\caption{Recorded results. The finite checks supplement the all-degree proof;
they do not prove it by extrapolation.}
\label{tab:checks}
\end{table}

\subsection{The independent endpoint-state calculation}

The verifier also computes an upper-bound relaxation by enumerating integer
pairs $(r,P)$ with $0\leq r,P\leq k$. In the total-degree calculation, it
retains the necessary condition
\[
d(P-r)<k-r,
\]
maximizes the relaxed count $k+P-r+2$ over these states, includes $k+1$ for a
shared endpoint winner, and caps the result by $2k+1$.
For the box calculation, it uses $d(P-r)<k$ and the elementary cap $2k+2$.
In the tested parameter range these independently computed upper bounds
match the certificate cardinalities exactly.

This is \emph{not} an exhaustive search over all coefficient vectors, all
supports, or all tropical polynomials. Some enumerated endpoint states need
not be realizable. Because they are only used to relax necessary conditions,
that does not invalidate the upper bound. Its all-degree validity is supplied
by Section~\ref{sec:upper}, not by the enumeration.

\subsection{Commands and files}

From the unpacked package directory, run
\begin{verbatim}
python3 code/verify.py --max-k 40
python3 code/select_area.py
pdflatex -interaction=nonstopmode -halt-on-error article.tex
pdflatex -interaction=nonstopmode -halt-on-error article.tex
\end{verbatim}
The second command audits the saved random-selection record without drawing
again. A deliberately requested new draw writes a separate output file.
The verifier should not be run with Python's \texttt{-O} flag, which disables
assertions; the entry point explicitly rejects that mode.

\texttt{data/verification.csv} contains one row per parameter--degree case,
including the least exact gap. The JSON summary records aggregate counts.
Separate JSON files contain the hand certificate and a degree-12 certificate
for $d=3$. The first 121 values of $H_3$ appear in
\texttt{data/hilbert\_d3.csv}. The source and proof audits are plain-text files
under \texttt{notes/}.

\section{Interpretation and limits}
\label{sec:limits}

\subsection{What has been established}

The counterexample is not based on an irrational slope, irrational exponent,
nonclosed set, limiting support, or ambiguous degree convention. It uses two
closed rational horizontal segments, nonnegative integer exponents, ordinary
total degree, and a finite family of rational coefficients at every $k$.
The coefficients are allowed to depend on $k$, as they must be when evaluating
a Hilbert function degree by degree. Within each fixed $k$, however, the same
coefficients serve every witness.

The exact formula supplies both a lower and an upper bound for every degree.
It therefore proves the degree limit as well as its nonintegrality. A single
finite certificate would not suffice to establish either asymptotic claim;
that is why the universal upper bound and uniform construction are both
included.

The geometric effect is visible even though the component directions never
change. Increasing the horizontal gap restricts how much the active
$x$-slope can reset between the components. The available reset is asymptotic
to $k/d$, and its integer rounding yields a genuine quasipolynomial.
Consequently, for this invariant, the relative placement of components
contains information not captured by their separate directions and lengths.

\subsection{What is left outside the claim}

The proof does not determine the Hilbert function of every rational
polyhedral set. It does not establish a general eventual-quasipolynomiality
theorem, a general rationality theorem for tropical degrees, or the existence
of higher-dimensional normalized limits. Nor does it settle the same
questions under a connectedness or balancing hypothesis.

The algebraic inequalities in the construction remain meaningful for a real
parameter $d>0$. We have deliberately stated the main geometric claim for
rational $d$, so it lies within the source's rational-polyhedral setting
without requiring a convention about irrational defining data. No assertion
about irrational tropical degree in that rational category is being made.

\subsection{A focused next mathematical question}

A natural continuation is whether every finite union of rational segments
has an eventually quasipolynomial strict-witness Hilbert function. The example
shows that polynomiality is too strong, but does not decide this replacement.
The endpoint argument suggests finite systems of integer inequalities whose
right-hand sides grow with $k$; proving that such endpoint data completely
characterize realizable supports would be the missing step, not a consequence
of the computations reported here.

\appendix
\section{Problem selection, exclusions, and source audit}
\label{app:audit}

\subsection{The random draw}

Before selecting a tropical problem, a fixed list of 96 mathematical areas
was constructed, and one call to Python's \texttt{secrets.randbelow(96)} was
made. The result was index 49 in zero-based indexing, or \textbf{50 in
one-based indexing: tropical geometry}. No redraw was made. The list and
the result are preserved in \texttt{data/area\_selection.json}; the full list
is printed in Appendix~\ref{app:areas}.

An operating-system random draw is not a deterministic sequence that another
reader should expect to reproduce. What is reproducible is the audit of the
saved draw, its selected list entry, and the subsequent calculations. The
packaged selection script defaults to this nonrandom audit; a new random draw
requires an explicit command and does not overwrite the retained record.

\subsection{The manifest exclusion check}

The supplied \texttt{manifest(1).tex} was read in full. It catalogues 71
research packages, and none of its listed problems concerns strict-witness
tropical Hilbert functions, tropical-degree integrality, or the two-segment
family studied here. Its own description distinguishes the claims in its
reports from independently checked proofs; for this task all listed problems
were excluded regardless of that distinction.

The SHA-256 digest of the supplied manifest is
\begin{center}
\small\ttfamily
729fb6e3aef3f30c7c5284c4296681fba72\\
b6d6af68211f178ad0bc3c5d15b41
\end{center}
The line break is typographical only: the digest has 64 hexadecimal digits.
The saved JSON record contains the uninterrupted digest.

\subsection{Version and priority qualification}

The arXiv abstract page retrieved on 20 September 2026 listed v2, dated
11 August 2025, as the latest version of 2404.06440. The v2 HTML and parsed
PDF text were consulted for the definitions and questions. PDF screenshot
requests failed with a fetch/cache error, so no claim is made that the
source's rendered PDF pages were visually inspected. The present article's
own compiled PDF was rendered and inspected separately.

Targeted searches for the paper title and for combinations of tropical
Hilbert function, counterexample, noninteger degree, and segments did not
locate a later resolution. Some searches were noisy or uninformative.
This supports the limited statement that the latest primary version retrieved
still poses the questions, not a guarantee that no unindexed preprint,
private proof, or independent concurrent work exists.

\subsection{Dependency audit}

The definitions and problem statements are the only mathematical inputs
needed from \cite{Grigoriev2025}. The degree computation does not invoke that
paper's general existence proof or its formulas for stars or individual
polyhedra. References~\cite{MaclaganRincon,Balanced} serve only to distinguish
the ideal-theoretic setting from the question answered here. Every bound and
certificate used in the result is proved within this article.

\section{Checklist for auditing the proof}
\label{app:proofaudit}

The following are the points at which an off-by-one error or a change of
convention would affect the conclusion.

\begin{enumerate}[leftmargin=1.6em]
\item The exponent region is $i+j\leq k$, not $i+j<k$ and not a square.
The separate box result has its own proof.
\item The support is one finite set with one common coefficient vector.
Every chosen monomial has a strict witness against every other chosen
monomial, not merely against monomials on its own segment.
\item The perturbation is inside an open set preserving all original strict
witnesses. Endpoint uniqueness is then obtained by avoiding finitely many
proper hyperplanes.
\item The count on a horizontal segment uses distinct integer slopes. Terms
that tie at an isolated point but never win strictly do not contribute an
additional active slope.
\item The two active sets may overlap. Their union is counted, and the
shared facing-winner case is separated before applying the strict comparison.
\item At distinct facing winners, adding the two strict inequalities gives
$d(P-r)<s-Q$, with this sign. The strict inequality implies
$P-r\leq\lceil k/d\rceil-1$ even when $k/d$ is integral.
\item The universal bound $2k+1$ uses the unique total-degree exponent
$(k,0)$ of $x$-slope $k$. It must not be copied to a box cutoff.
\item In the lower bound, $m<k$ makes the two rows of exponent vectors
disjoint. The slack $\sigma$ is positive for both $d\leq1$ and $d\geq1$.
\item The witnesses lie in the interiors of the required segments. The
cross-group gaps are explicitly positive, including at the transition $d=1$
and when $k/d$ is an integer.
\item The case $k=0$ is handled separately before dividing by $k$ in the
certificate parameters. Limits and nonpolynomiality are then deduced from
the all-degree formula, not from the finite verification range.
\end{enumerate}

\section{The complete random-selection list}
\label{app:areas}

The ordering below is the ordering supplied to the random draw. Both columns
use the original one-based indices; they do not define separate draws.

\small
\begin{longtable}{@{}r p{6.55cm} r p{6.55cm}@{}}
\toprule
No.&Area&No.&Area\\
\midrule
\endfirsthead
\toprule
No.&Area&No.&Area\\
\midrule
\endhead
1 & Additive combinatorics & 49 & Transcendental number theory \\
2 & Algebraic combinatorics & 50 & \textbf{Tropical geometry} \\
3 & Analytic number theory & 51 & Word combinatorics \\
4 & Arithmetic geometry & 52 & Zero-sum theory \\
5 & Automata and formal languages & 53 & Association schemes \\
6 & Banach space geometry & 54 & Coding theory \\
7 & Combinatorial commutative algebra & 55 & Discrete optimization \\
8 & Combinatorial design theory & 56 & Extremal set theory \\
9 & Combinatorial game theory & 57 & Frame theory \\
10 & Computable analysis & 58 & Graph colorings \\
11 & Computational geometry & 59 & Incidence geometry \\
12 & Convex geometry & 60 & Infinite-dimensional holomorphy \\
13 & Discrete geometry & 61 & Matrix inequalities \\
14 & Dynamical systems & 62 & Moment problems \\
15 & Enumerative geometry & 63 & Orthogonal polynomials \\
16 & Ergodic theory & 64 & Positive definite functions \\
17 & Extremal graph theory & 65 & Rational approximation \\
18 & Finite fields & 66 & Riemannian geometry \\
19 & Finite group theory & 67 & Simplicial complexes \\
20 & Functional equations & 68 & Tensor rank \\
21 & General topology & 69 & Topological dynamics \\
22 & Geometric group theory & 70 & Universal algebra \\
23 & Graph polynomials & 71 & Ultrametric analysis \\
24 & Hypergraph theory & 72 & Valuation theory \\
25 & Integer partitions & 73 & Boolean function theory \\
26 & Integral inequalities & 74 & Chip-firing and sandpiles \\
27 & Knot invariants & 75 & Discrete probability \\
28 & Lattice theory & 76 & Fractional calculus \\
29 & Linear preserver problems & 77 & Group actions on trees \\
30 & Matroid theory & 78 & Hyperplane arrangements \\
31 & Metric geometry & 79 & Information inequalities \\
32 & Noncommutative algebra & 80 & Modular forms \\
33 & Numerical semigroups & 81 & Monomial ideals \\
34 & Operator inequalities & 82 & Optimal transport \\
35 & Ordered algebraic structures & 83 & Partial orders \\
36 & Permutation patterns & 84 & Real-rooted polynomials \\
37 & Polynomial dynamics & 85 & Rewriting systems \\
38 & Probabilistic combinatorics & 86 & Schubert calculus \\
39 & Quasigroups and loops & 87 & Set-theoretic topology \\
40 & Ramsey theory & 88 & Special functions \\
41 & Real algebraic geometry & 89 & String-rewriting and monoids \\
42 & Recurrence sequences & 90 & Sum-product phenomena \\
43 & Riordan arrays & 91 & Topological semigroups \\
44 & Semigroup theory & 92 & Visibility in graphs \\
45 & Spectral graph theory & 93 & Weighted lattice paths \\
46 & Symbolic dynamics & 94 & Well-quasi-orders \\
47 & Tiling theory & 95 & Zeta functions of graphs \\
48 & Topological graph theory & 96 & Numerical linear algebra \\
\bottomrule
\end{longtable}
\normalsize

\begin{thebibliography}{9}
\bibitem{Grigoriev2025}
D. Grigoriev,
\emph{Degree of 1-dimensional tropical prevariety},
arXiv:2404.06440v2, 11 August 2025.
Definitions 1.1--1.2; introduction, eventual-polynomiality conjecture and
Question 2.
\url{https://arxiv.org/abs/2404.06440v2}.
Versioned full text:
\url{https://arxiv.org/html/2404.06440v2}.

\bibitem{Grigoriev2024}
D. Grigoriev,
\emph{A tropical version of Hilbert function},
arXiv:2404.06440v1, 9 April 2024.
Cited only to identify the earlier title and version.
\url{https://arxiv.org/abs/2404.06440v1}.

\bibitem{MaclaganRincon}
D. Maclagan and F. Rinc\'on,
\emph{Tropical ideals},
Compositio Mathematica \textbf{154} (2018), 640--670.
\href{https://doi.org/10.1112/S0010437X17008004}{doi:10.1112/S0010437X17008004}.
\url{https://arxiv.org/abs/1609.03838}.

\bibitem{Balanced}
D. Maclagan and F. Rinc\'on,
\emph{Varieties of Tropical Ideals are Balanced},
arXiv:2009.14557, 30 September 2020.
\url{https://arxiv.org/abs/2009.14557}.
\end{thebibliography}

\end{document}
